% Generated draft: requires mathematical and source review.
\documentclass[11pt]{article}
\usepackage[margin=25mm]{geometry}
\usepackage{fontspec}
\setmainfont{Latin Modern Roman}
\usepackage{amsmath,amssymb,mathtools,unicode-math}
\setmathfont{Latin Modern Math}
\usepackage{enumitem,needspace,xurl,hyperref}
\setlength{\parindent}{0pt}
\setlength{\parskip}{5pt}
\setlength{\emergencystretch}{3em}
\newcommand{\N}{\mathbb{N}}
\newcommand{\Z}{\mathbb{Z}}
\newcommand{\Q}{\mathbb{Q}}
\newcommand{\R}{\mathbb{R}}
\newcommand{\C}{\mathbb{C}}
\newcommand{\F}{\mathbb{F}}
\newcommand{\A}{\mathbb{A}}
\newcommand{\PP}{\mathbb{P}}
\newcommand{\E}{\mathbb{E}}
\newcommand{\eps}{\varepsilon}
\newcommand{\dd}{\,\mathrm{d}}
\begin{document}
\section*{UnsolvedMath ID: 700455}
\textbf{Source problem code:} catalogue entry
\subsection{Definitions and mathematical statement}
Let \(m,n\in\Z_{>0}\), let \(p\) be a prime, and put \(R_m=\Z/p^m\Z\). Write \(\rho:R_m\to\F_p\) for reduction modulo \(p\), also applied entrywise to matrices. Define
\[
U_n(\F_p)=\{(u_{ij})\in\operatorname{GL}_n(\F_p):u_{ii}=1\text{ for all }i,\ u_{ij}=0\text{ for }i>j\}.
\]
The supplied statement describes \(P_n(R_m)\) by requiring \(a_{ii}\equiv1\) for every \(i\), without specifying the modulus, and \(a_{ij}\equiv0\pmod p\) for \(i>j\). If the diagonal congruence is also modulo \(p\), this gives the finite \(p\)-group
\[
P_n(R_m)=\{A=(a_{ij})\in\operatorname{GL}_n(R_m):\rho(A)\in U_n(\F_p)\}.
\]
There is no further restriction on entries above the diagonal. This interpretation of the diagonal condition requires confirmation: imposing \(a_{ii}=1\) in \(R_m\) instead does not in general give a group under matrix multiplication. The formulation below is conditional on the modulo-\(p\) interpretation and does not silently resolve the source omission.

For a group \(H\), write \(H'=[H,H]\) for its commutator subgroup. A finite \(p\)-group \(G\) is called \emph{regular} if, for every \(x,y\in G\), there exist an integer \(r\geq0\) and elements \(c_1,\ldots,c_r\in\langle x,y\rangle'\) such that
\[
(xy)^p=x^py^p c_1^p\cdots c_r^p,
\]
where the product is the identity when \(r=0\).

\textbf{Problem (B. A. F. Wehrfritz, 1982).} For every prime \(p\), determine exactly which pairs \((m,n)\in\Z_{>0}^2\) make \(P_n(\Z/p^m\Z)\) regular, subject to confirmation of the diagonal congruence above.

\textbf{Supplied editorial qualifications.} The editors' comment dated 2001 reports that regularity holds when \(mn<p\), citing Yu. I. Merzlyakov [3]. It reports that A. V. Yagzhev [4] completely treated \(m=1\), while explicitly warning that that paper is erroneous for \(m>1\). It also reports that S. G. Kolesnikov [5] completed the case \(m=2\). The editors' comment dated 2009 reports that Kolesnikov [6] proved regularity for every positive integer \(m\) whenever \(n^2<p\). These are supplied historical assessments; the complete classifications for \(m=1\) and \(m=2\) are not stated in the record. The local registry labels the problem open; that label is not an independently verified literature assessment.
\subsection{Short English statement}
For each prime \(p\), classify the positive integers \(m,n\) for which the indicated matrix group over \(\Z/p^m\Z\) is a regular finite \(p\)-group. The source omits the modulus in its diagonal condition, so this draft conditionally reads that condition modulo \(p\). Preserve the supplied results and warnings for special cases; the general open label has not been independently verified.
\subsection{Sources}
\begin{itemize}
\item[{[1]}] Local UnsolvedMath registry, supplied complete record with permanent ID \(700455\), exact source problem code \texttt{catalogue entry}, title \emph{Kourovka Notebook Problem 8.3}. Supplied status: open. Upstream revision: \texttt{d371ee917ac9647fbb06cd6fc529deec6f3dda82}; record updated \texttt{2026-10-16T14:00:00Z}.
\item[{[2]}] \emph{The Kourovka Notebook (Unsolved Problems in Group Theory)}, Problem 8.3, proposed by B. A. F. Wehrfritz in the 8th issue (1982). The supplied record identifies the 21st-issue edition (2026), \texttt{arXiv:1401.0300v48}. The collection description credits editors E. I. Khukhro and V. D. Mazurov. Editorial comments dated 2001 and 2009 are reproduced as supplied historical assessments, not independently verified findings.
\item[{[3]}] Yu. I. Merzlyakov, \emph{Algebra i Logika} \textbf{3}, no. 4 (1964), 49--59 (Russian). The supplied 2001 editorial comment attributes regularity under \(mn<p\) to this reference; no article title is supplied.
\item[{[4]}] A. V. Yagzhev, \emph{Math. Notes} \textbf{56}, nos. 5--6 (1995), 1283--1290. The supplied 2001 editorial comment says that this paper completely treats \(m=1\), but is erroneous for \(m>1\); no article title is supplied.
\item[{[5]}] S. G. Kolesnikov, \emph{Issledovaniya in analysis and algebra}, no. 3, Tomsk Univ., Tomsk (2001), 117--125 (Russian), with the publication name reproduced as supplied. The supplied 2001 editorial comment attributes completion of \(m=2\) to this reference; no article title is supplied.
\item[{[6]}] S. G. Kolesnikov, \emph{Siberian Math. J.} \textbf{47}, no. 6 (2006), 1054--1059. The supplied 2009 editorial comment attributes regularity for every \(m\) when \(n^2<p\) to this reference; no article title is supplied.
\end{itemize}
\end{document}
